\documentclass{beamer}
\usetheme{Madrid}
\usepackage[spanish]{babel}
\usepackage[utf8]{inputenc}
\usepackage{graphicx}
\usepackage{hyperref}
\usepackage{tikz}
\usetikzlibrary{shapes.geometric, arrows.meta, positioning}

\title{Finanzas P\'ublicas Colombianas}
\author{Oliver Pardo y Andres Garzón}
\institute{Pontificia Universidad Javeriana}
\date{2025}

\begin{document}

\begin{frame}
  \titlepage
\end{frame}

\begin{frame}
  \frametitle{Datos de Contacto}
  \centering
  \Large \textit{Oliver Pardo} \\ \vspace{0.3cm}
  \Large \textit{PhD en Economía, LSE.} \\ \vspace{0.3cm}
  \Large \textit{pardoo@javeriana.edu.co} \\ \vspace{0.3cm}
  \large{Profesor Asociado FCEA PUJ} \\ \vspace{0.3cm}
  \large{Ex Director Política Macroeconómica MHCP} \\ \vspace{0.3cm}
  \large{Ex Director Observatorio Fiscal PUJ} \\ 
\end{frame}


\begin{frame}
  \frametitle{Introducci\'on}
  El curso ofrece una mirada al funcionamiento y los retos del sector p\'ublico en Colombia. Estudia c\'omo se financia el Estado, c\'omo se asignan los recursos p\'ublicos y sus implicaciones para el desarrollo econ\'omico, la equidad y la sostenibilidad fiscal.
\end{frame}

\begin{frame}
  \frametitle{Objetivos del Curso}
  \begin{itemize}
    \item Comprender la estructura del sector p\'ublico colombiano.
    \item Analizar ingresos y gastos p\'ublicos.
    \item Evaluar la sostenibilidad de las finanzas p\'ublicas.
    \item Interpretar documentos oficiales clave.
    \item Formular propuestas de pol\'itica fiscal.
    \item Desarrollar habilidades en manejo de datos fiscales.
  \end{itemize}
\end{frame}

\begin{frame}
  \frametitle{Metodolog\'ia y Evaluaci\'on}
  \textbf{Metodolog\'ia:} Clase magistral (3h) + taller aplicado (1h).\\
  \textbf{Evaluaci\'on:}
  \begin{itemize}
    \item Quizzes: 33.3\%
    \item Presentaci\'on en grupo: 33.3\%
    \item Taller: 33.3\%
  \end{itemize}
\end{frame}


\begin{frame}
  \frametitle{Contenido del Curso}
  M\'odulos:
  \begin{itemize}
    \item M0: Presentaci\'on del curso
    \item M1: Sector p\'ublico colombiano
    \item M2: Gobierno Nacional Central
    \item M3: Deuda Pública
    \item M4: Regla fiscal y ciclo presupuestal
    \item M5: Seguridad social en pensiones
    \item M6: Seguridad social en salud
    \item M7: Finanzas p\'ublicas territoriales
    \item M8: Comparaci\'on internacional
    \item M9: Presentaciones en grupo
  \end{itemize}
\end{frame}

\begin{frame}
  \frametitle{Componente Aplicado: Manejo de Datos}
  \begin{itemize}
    \item Bases de datos oficiales (PGN, SGP, GEIH, FUT, Terridata, SIIF).
    \item Talleres pr\'acticos con datos reales.
    \item Visualizaci\'on e interpretaci\'on de indicadores fiscales.
    \item Enlace directo entre teor\'ia econ\'omica y evidencia emp\'irica.
  \end{itemize}
\end{frame}

\begin{frame}
  \frametitle{Funciones del Estado y su Base Constitucional}
  \begin{itemize}
    \item \textbf{Legislativa:} Crear leyes (Art. 1, 3)
    \item \textbf{Ejecutiva:} Administrar y ejecutar leyes (Art. 113, 209)
    \item \textbf{Judicial:} Aplicar justicia (Art. 229, 228)
    \item \textbf{Defensa y seguridad:} Proteger soberan\'ia y orden (Art. 2)
    \item \textbf{Econ\'omica:} Dirigir la econom\'ia y corregir desigualdades (Art. 1, 334)
    \item \textbf{Social:} Garantizar derechos b\'asicos (Art. 1, 2, 366)
    \item \textbf{Fiscal:} Recaudar y asignar recursos (Art. 338, 95)
    \item \textbf{Reguladora:} Supervisar mercados y servicios (Art. 333, 365)
  \end{itemize}
\end{frame}

\begin{frame}
  \frametitle{Art\'iculo 1 - CP de 1991}
  Colombia es un Estado social de derecho, organizado en forma de Rep\'ublica unitaria, descentralizada, con autonom\'ia de sus entidades territoriales, democr\'atica, participativa y pluralista, fundada en el respeto de la dignidad humana, en el trabajo y la solidaridad de las personas que la integran y en la prevalencia del inter\'es general.
\end{frame}

\begin{frame}
  \frametitle{Art\'iculo 2 - CP de 1991}
  Son fines esenciales del Estado: servir a la comunidad, promover la prosperidad general y garantizar la efectividad de los principios, derechos y deberes consagrados en la Constituci\'on; facilitar la participaci\'on de todos en las decisiones que los afectan y en la vida econ\'omica, pol\'itica, administrativa y cultural de la Naci\'on; defender la independencia nacional, mantener la integridad territorial y asegurar la convivencia pac\'ifica y la vigencia de un orden justo.
\end{frame}

\begin{frame}
  \frametitle{Art\'iculo 3 - CP de 1991}
  La soberan\'ia reside exclusivamente en el pueblo, del cual emana el poder p\'ublico. El pueblo la ejerce en forma directa o por medio de sus representantes, en los t\'erminos que la Constituci\'on establece.
\end{frame}

\begin{frame}
  \frametitle{Art\'iculo 95 - CP de 1991}
  La constituci\'on establece como deber de la persona y del ciudadano contribuir al financiamiento de los gastos e inversiones del Estado dentro de conceptos de justicia y equidad.
\end{frame}

\begin{frame}
  \frametitle{Art\'iculo 113 - CP de 1991}
  Son ramas del poder p\'ublico la legislativa, la ejecutiva y la judicial. Adem\'as de los \'organos que las integran, existen otros, aut\'onomos e independientes, para el cumplimiento de las funciones del Estado. Los diferentes \'organos del Estado tienen funciones separadas pero colaboran arm\'onicamente para la realizaci\'on de sus fines.
\end{frame}

\begin{frame}
  \frametitle{Art\'iculo 209 - CP de 1991}
  La funci\'on administrativa est\'a al servicio de los intereses generales y se desarrolla con fundamento en los principios de igualdad, moralidad, eficacia, econom\'ia, celeridad, imparcialidad y publicidad, mediante la descentralizaci\'on, la delegaci\'on y la desconcentraci\'on de funciones.
\end{frame}

\begin{frame}
  \frametitle{Art\'iculo 228 - CP de 1991}
  La administraci\'on de justicia es funci\'on p\'ublica. Sus decisiones son independientes. Las actuaciones ser\'an p\'ublicas y permanentes, con las excepciones que establezca la ley, y los t\'erminos procesales se observar\'an con diligencia. Su funcionamiento ser\'a desconcentrado y aut\'onomo.
\end{frame}

\begin{frame}
  \frametitle{Art\'iculo 229 - CP de 1991}
  Se garantiza el derecho de toda persona para acceder a la administraci\'on de justicia. La ley indicar\'a en qu\'e casos puede hacerlo sin la representaci\'on de abogado.
\end{frame}

\begin{frame}
  \frametitle{Art\'iculo 333 - CP de 1991}
  La actividad econ\'omica y la iniciativa privada son libres, dentro de los l\'imites del bien com\'un. La empresa, como base del desarrollo, tiene una funci\'on social que implica obligaciones. El Estado, por mandato de la ley, intervendr\'a por razones de utilidad p\'ublica o de inter\'es social.
\end{frame}

\begin{frame}
  \frametitle{Art\'iculo 334 - CP de 1991}
  El Estado intervendr\'a por mandato de la ley en la explotaci\'on de los recursos naturales, en el uso del suelo, en la producci\'on, distribuci\'on, utilizaci\'on y consumo de los bienes, y en los servicios p\'ublicos y privados, para racionalizar la econom\'ia con el fin de conseguir el mejoramiento de la calidad de vida.
\end{frame}

\begin{frame}
  \frametitle{Art\'iculo 338 - CP de 1991}
  En tiempo de paz, solamente el Congreso, las asambleas departamentales y los concejos distritales y municipales podr\'an imponer contribuciones fiscales o parafiscales. La ley, las ordenanzas y los acuerdos deben fijar directamente los sujetos activos y pasivos, los hechos y las bases gravables, y las tarifas de los impuestos.
\end{frame}

\begin{frame}
  \frametitle{Art\'iculo 365 - CP de 1991}
  Los servicios p\'ublicos son inherentes a la finalidad social del Estado. Es deber del Estado asegurar su prestaci\'on eficiente a todos los habitantes del territorio nacional. Podr\'an ser prestados por el Estado, directa o indirectamente, por comunidades organizadas, o por particulares.
\end{frame}

\begin{frame}
  \frametitle{Art\'iculo 366 - CP de 1991}
  El bienestar general y el mejoramiento de la calidad de vida de la poblaci\'on son finalidades sociales del Estado. Ser\'a objetivo fundamental de su actividad la soluci\'on de las necesidades insatisfechas de salud, de educaci\'on, de saneamiento ambiental y de agua potable.
\end{frame}


\begin{frame}
  \frametitle{Funciones del Estado y su Base Constitucional}
  \begin{itemize}
    \item \textbf{Legislativa:} Crear leyes (Art. 1, 3)
    \item \textbf{Ejecutiva:} Administrar y ejecutar leyes (Art. 113, 209)
    \item \textbf{Judicial:} Aplicar justicia (Art. 229, 228)
    \item \textbf{Defensa y seguridad:} Proteger soberan\'ia y orden (Art. 2)
    \item \textbf{Econ\'omica:} Dirigir la econom\'ia y corregir desigualdades (Art. 1, 334)
    \item \textbf{Social:} Garantizar derechos b\'asicos (Art. 1, 2, 366)
    \item \textbf{Fiscal:} Recaudar y asignar recursos (Art. 338, 95)
    \item \textbf{Reguladora:} Supervisar mercados y servicios (Art. 333, 365)
  \end{itemize}
\end{frame}

\begin{frame}{Desagregaci\'on del Sector P\'ublico Consolidado}
\centering
\begin{tikzpicture}[node distance=0.5cm and 0.5cm, every node/.style={align=center},
  box/.style={rectangle, draw, rounded corners, text width=3cm, minimum height=0.7cm},
  level/.style={rectangle, draw, text width=3cm, minimum height=0.7cm}]

\node[box, fill=blue!10] (spc) {\tiny{Sector P\'ublico Consolidado}};

\node[level, below left=of spc] (spf) {\tiny{Sector Público Financiero}};
\node[level, fill=green!10, below =of spc] (spnf) {\tiny{Sector Público No Financiero}};


\node[level, fill=green!10, below=of spnf] (gg) {\tiny{Gobierno General}};
\node[level, below left =of spnf] (eepp) {\tiny{Empresas Públicas}};


\node[level, fill=red!10, below left=of gg] (gnc) {\tiny{Gobierno Nacional Central}};
\node[level, fill=red!10, below =of gg] (et) {\tiny{Entidades Territoriales}};
\node[level, fill=red!10, below right =of gg] (ss) {\tiny{Seguridad Social}};

%\node[level, below left=of ss] (salud) {\tiny{Salud}};
%\node[level, below right =of ss] (pensiones) {\tiny{Pensiones}};

\node[box, fill=blue!10] (spc) {Sector P\'ublico Consolidado};

% Conexiones directas desde SPC
\draw[->] (spc) -- (spf);
\draw[->] (spc) -- (spnf);

\draw[->] (spnf) -- (gg);
\draw[->] (spnf) -- (eepp);

\draw[->] (gg) -- (gnc);
\draw[->] (gg) -- (et);
\draw[->] (gg) -- (ss);
\end{tikzpicture}
\end{frame}

\begin{frame}
  \frametitle{Gracias}
\end{frame}

\end{document}
